% #################################################################
% JGME Manuscript + Cover Letter Template
% Journal of Graduate Medical Education
% Generated 2026
%
% This template separates:
%   1. Cover letter
%   2. Title page
%   3. Manuscript
%   4. References
%   5. Tables/Figures/Boxes
%   6. Supplemental material checklist
%
% IMPORTANT:
% - Choose ONE manuscript category and obey its word/figure limits.
% - JGME research submissions use structured abstracts/manuscripts.
% - Perspectives and On Teaching do not require an abstract.
% - The journal checklist requires double-spaced, numbered pages.
% - References are numbered consecutively in order of first citation.
% - Research manuscripts require study year, response-rate reporting
%   when applicable, IRB/ethics statement, limitations, and survey
%   supplement when applicable.
% - AI-related submissions should clearly disclose AI use and retain
%   human responsibility for final content.
% #################################################################

\documentclass[12pt]{article}

\usepackage[margin=1in]{geometry}
\usepackage{setspace}
\usepackage{lineno}
\usepackage{enumitem}
\usepackage{booktabs}
\usepackage{longtable}
\usepackage{graphicx}
\usepackage{caption}
\usepackage{hyperref}
\usepackage{url}
\usepackage{microtype}
\usepackage{titlesec}
\usepackage{fancyhdr}

\hypersetup{
  hidelinks,
  pdfauthor={AUTHOR NAME},
  pdftitle={MANUSCRIPT TITLE}
}

\doublespacing
\linenumbers

\pagestyle{fancy}
\fancyhf{}
\fancyfoot[C]{\thepage}

\setlength{\parindent}{0.5in}
\setlength{\parskip}{0pt}

% ------------------------------------------------------------------
% USER VARIABLES
% ------------------------------------------------------------------

\newcommand{\ManuscriptTitle}{FULL MANUSCRIPT TITLE}
\newcommand{\ShortTitle}{SHORT RUNNING TITLE}
\newcommand{\ManuscriptCategory}{Perspective} % Change as needed.
\newcommand{\CorrespondingAuthor}{FULL NAME, DEGREE(S)}
\newcommand{\CorrespondingInstitution}{INSTITUTION}
\newcommand{\CorrespondingAddress}{BUSINESS ADDRESS}
\newcommand{\CorrespondingPhone}{PHONE}
\newcommand{\CorrespondingEmail}{EMAIL}
\newcommand{\ManuscriptWordCount}{0000}
\newcommand{\AbstractWordCount}{000}

% JGME category limits:
%
% Original Research:
%   Main text <= 2,500 words
%   Qualitative research <= 3,500 words
%   Abstract <= 250 words
%   Tables/Figures/Boxes <= 5
%
% Educational Innovation:
%   Main text <= 2,000 words
%   Abstract <= 250 words
%   Tables/Figures/Boxes <= 5
%
% Review:
%   Main text <= 3,000 words
%   Abstract <= 250 words
%   Tables/Figures/Boxes <= 5
%
% Brief Report:
%   Main text <= 1,200 words
%   Abstract <= 250 words
%   Tables/Figures/Boxes <= 2
%
% Perspective:
%   Main text <= 1,200 words
%   No abstract
%   Tables/Figures/Boxes <= 2
%
% On Teaching:
%   Main text <= 1,200 words
%   No abstract
%
% To the Editor:
%   Main text <= 500 words
%   No abstract
%   Tables/Figures/Boxes <= 2

\begin{document}

% ==================================================================
% COVER LETTER
% ==================================================================

\begin{flushleft}
	\today

	Editor-in-Chief\\
	Journal of Graduate Medical Education
\end{flushleft}

\vspace{1em}

Dear Editor:

Please consider our manuscript, ``\ManuscriptTitle,'' for publication in the
\textit{Journal of Graduate Medical Education} as a \ManuscriptCategory.

% ------------------------------------------------------------------
% Required/strongly recommended cover-letter content
% ------------------------------------------------------------------

This manuscript addresses the following graduate medical education problem:
[ONE TO THREE SENTENCES EXPLAINING THE GME PROBLEM, POPULATION, OR NEED].

The principal contribution of this work is:
[ONE TO THREE SENTENCES EXPLAINING WHAT IS NEW, USEFUL, OR ACTIONABLE].

This manuscript is relevant to JGME readers because:
[ONE TO THREE SENTENCES LINKING THE WORK TO RESIDENCY/FELLOWSHIP EDUCATION,
PROGRAM LEADERSHIP, FACULTY DEVELOPMENT, ASSESSMENT, SUPERVISION, OR GME
SCHOLARSHIP].

% For AI-related manuscripts, consider retaining the next paragraph.
% Delete if not applicable.
This manuscript addresses artificial intelligence in graduate medical
education. We have described where and how AI tools were used in the study,
educational intervention, analysis, and/or manuscript preparation; all
authors reviewed and take responsibility for the final content. Protected
health information was not entered into unapproved consumer AI systems.
	[MODIFY THIS LANGUAGE TO MATCH WHAT ACTUALLY OCCURRED.]

% Originality / duplicate-publication statement.
This manuscript is original, has not been published previously, and is not
under consideration by another journal. [MODIFY IF THERE IS A PREPRINT,
		ABSTRACT, CONFERENCE POSTER, OR RELATED PUBLICATION; DISCLOSE IT HERE.]

% Authorship.
All listed authors meet applicable authorship criteria, approved the submitted
version, and agree to be accountable for the work.

	% Ethics.
	[SELECT/EDIT ONE:]
The study received approval from [IRB/ETHICS BODY], protocol [NUMBER].
OR
The study was determined exempt by [IRB/ETHICS BODY], determination [NUMBER].
OR
This manuscript does not report human-subjects research requiring IRB review.

% Conflicts/funding.
Conflicts of interest: [NONE / DESCRIBE].
Funding/support: [NONE / DESCRIBE].

% Optional word-limit explanation.
% JGME has specifically noted that an over-limit manuscript should explain
% the reason in the cover letter.
[IF THE MANUSCRIPT EXCEEDS THE CATEGORY WORD LIMIT, EXPLAIN AND JUSTIFY THE
EXCESS HERE. OTHERWISE DELETE THIS PARAGRAPH.]

% Optional AI Teaching Rounds note.
[OPTIONAL: We believe this manuscript may be relevant to JGME's AI Teaching
Rounds conversation because ... . We understand that JGME may classify the
submission under its standard article categories.]

Thank you for considering this manuscript.

Sincerely,

\CorrespondingAuthor\\
\CorrespondingInstitution\\
\CorrespondingAddress\\
\CorrespondingPhone\\
\CorrespondingEmail

\clearpage

% ==================================================================
% TITLE PAGE
% ==================================================================

\begin{center}
	{\Large\bfseries \ManuscriptTitle\par}

	\vspace{1em}

	\textbf{Manuscript category:} \ManuscriptCategory

	\vspace{1em}

	% List every author, degree, affiliation, and role.
	AUTHOR 1, DEGREE(S)$^{1}$\\
	Role: [eg, resident, fellow, program director, associate professor]\\[0.5em]

	AUTHOR 2, DEGREE(S)$^{2}$\\
	Role: [ROLE]\\[0.5em]

	AUTHOR 3, DEGREE(S)$^{1,3}$\\
	Role: [ROLE]
\end{center}

\vspace{1em}

\textbf{Affiliations}

$^{1}$ Department/Program, Institution, City, State/Province, Country

$^{2}$ Department/Program, Institution, City, State/Province, Country

$^{3}$ Department/Program, Institution, City, State/Province, Country

\vspace{1em}

\textbf{Corresponding author}

\CorrespondingAuthor\\
\CorrespondingInstitution\\
\CorrespondingAddress\\
Phone: \CorrespondingPhone\\
Email: \CorrespondingEmail

\vspace{1em}

\textbf{Word counts}

Main manuscript: \ManuscriptWordCount\ words

Abstract: \AbstractWordCount\ words
% For Perspective, On Teaching, and To the Editor use:
% Abstract: Not applicable

\vspace{1em}

\textbf{Abbreviations}

AI, artificial intelligence; GME, graduate medical education; [ADD OTHERS].

\vspace{1em}

\textbf{Funding/support}

[None / funding source, grant number, sponsor role.]

\vspace{1em}

\textbf{Conflicts of interest}

[None / describe.]

\vspace{1em}

\textbf{Prior presentation}

[None / meeting name, location, date.]

\vspace{1em}

\textbf{Data availability}

[Describe availability, repository, restrictions, or “Not applicable.”]

\vspace{1em}

\textbf{AI disclosure}

[If applicable: identify the AI system/tool, model/version when known,
	date/access context where relevant, what task it performed, what human review
	occurred, and confirm that authors remain responsible for final content.
	Delete if no AI was used and journal policy does not require a negative
	statement.]

\clearpage

% ==================================================================
% ABSTRACT
% ==================================================================

% DELETE THE ENTIRE ABSTRACT SECTION FOR:
%   Perspective
%   On Teaching
%   To the Editor
%
% USE A STRUCTURED ABSTRACT FOR:
%   Original Research
%   Educational Innovation
%   Brief Report
%
% JGME limit for these categories: <=250 words.

\section*{Abstract}

\textbf{Background:}
[What is already known? What GME problem or gap motivates the work?]

\textbf{Objective:}
[State the primary objective/question explicitly.]

\textbf{Methods:}
[Design, setting, participants, intervention/exposure, measures, analysis.
	INCLUDE THE YEAR(S) THE STUDY WAS CONDUCTED.]

\textbf{Results:}
[Primary findings with numerical data where appropriate. IF RESPONSE RATE IS
	APPLICABLE, INCLUDE TOTAL POSSIBLE N, ACTUAL N, AND CALCULATED PERCENTAGE.]

\textbf{Conclusions:}
[Answer the objective without overstating findings.]

\clearpage

% ==================================================================
% MAIN MANUSCRIPT: RESEARCH / EDUCATIONAL INNOVATION / BRIEF REPORT
% ==================================================================
%
% For research-style submissions use the following sections.
% For Perspective/On Teaching, delete this structure and use the
% alternative narrative structure below.

\section{Introduction}

 % Recommended logic:
 % 1. What is known?
 % 2. What remains unknown/problematic?
 % 3. Why does the gap matter to GME?
 % 4. What did this study/intervention aim to determine?

 [INTRODUCTION TEXT.]

\section{Methods}

\subsection{Study Design and Setting}

[Design, institution/program, specialty, geographic context, dates/year.]

\subsection{Participants}

[Eligibility, recruitment, sampling, number eligible, number participating.]

\subsection{Intervention or Exposure}

[Describe sufficiently for replication.]

% AI-specific optional subsection:
\subsection{Artificial Intelligence System and Workflow}

[If applicable: tool/model/provider, version or date, institutional approval,
	input/output workflow, prompts/protocol, human verification, PHI handling,
	access conditions, and reproducibility considerations.]

\subsection{Outcomes and Measures}

[Primary/secondary outcomes, instruments, validity evidence, operational
	definitions.]

\subsection{Analysis}

[Statistical and/or qualitative analysis; software and version where relevant.]

\subsection{Ethics}

% JGME requires the IRB/ethics approval or exemption statement at the END
% of the Methods section for research manuscripts.
[Institutional review board / ethical review approval, exemption, or other
	determination, including institution and identifier when applicable.]

\section{Results}

 % If response rate applies, report:
 % total possible participants, actual participants, calculated percentage.

 [RESULTS TEXT.]

\section{Discussion}

\subsection{Principal Findings}

[Interpret the findings; do not simply repeat Results.]

\subsection{Relationship to Prior Literature}

[Compare/contrast with existing GME literature.]

\subsection{Implications for Graduate Medical Education}

[What should educators/programs/researchers do differently?]

\subsection{Limitations}

% JGME explicitly requires a limitations statement for research submissions.
[State limitations clearly and proportionately.]

\section{Conclusions}

 [Concise answer to the study objective supported by the data.]

% ==================================================================
% ALTERNATIVE NARRATIVE STRUCTURE:
% PERSPECTIVE / AI TEACHING-ROUNDS-STYLE ARTICLE
% ==================================================================
%
% Delete the research sections above and use something like this for a
% Perspective. JGME Perspectives: <=1,200 words; no abstract; <=2
% tables/figures/boxes.
%
% \section{The Problem}
% [Concrete GME problem/case.]
%
% \section{Why It Matters Now}
% [Evidence and GME context.]
%
% \section{Principles}
% [2--4 practical principles.]
%
% \section{Practical Approach}
% [Actionable workflow.]
%
% \section{Risks and Guardrails}
% [Privacy, bias, supervision, deskilling, accountability, etc.]
%
% \section{Try This}
% [Small intervention readers can implement immediately.]
%
% \section{Next Steps}
% [What programs/researchers should evaluate next.]

% ==================================================================
% OPTIONAL PRACTICE BOX
% ==================================================================

\section*{Box 1. Practical Implementation Checklist}

\begin{itemize}[leftmargin=*]
	\item Define the educational problem before choosing an AI/tool solution.
	\item Establish what uses require disclosure.
	\item Identify tasks that must remain under direct human judgment.
	\item Use institution-approved systems for protected/sensitive information.
	\item Require human verification of generated output.
	\item Monitor effects on trainee reasoning and independent skill development.
	\item Evaluate equity of access to required tools.
	\item Start with a small pilot, measure outcomes, revise, then scale.
\end{itemize}

% ==================================================================
% REFERENCES
% ==================================================================

\section*{References}

% JGME requires references numbered consecutively in order of first
% appearance in text AND graphics.
%
% Use the journal's biomedical/NLM-like format.
%
% Example:
% 1. Author AA, Author BB, Author CC, et al. Article title.
%    J Grad Med Educ. 2026;18(1):1-6. doi:xx.xxxx/xxxxx

\begin{enumerate}
	\item Author AA, Author BB. Article title. \textit{Journal Name}.
	      YEAR;VOLUME(ISSUE):PAGES. doi:DOI
	\item Author AA, Author BB, Author CC, et al. Article title.
	      \textit{Journal Name}. YEAR;VOLUME(ISSUE):PAGES. doi:DOI
\end{enumerate}

% ==================================================================
% TABLES / FIGURES / BOXES
% ==================================================================

\clearpage
\section*{Tables, Figures, and Boxes}

% Every table, figure, box, and supplement must be cited in the manuscript.
% JGME checklist:
% - legends required;
% - footnotes use superscripted letters;
% - figures provided in original format;
% - grayscale;
% - minimum 200 dpi.

\begin{table}[ht]
	\centering
	\caption{Example Table Title}
	\begin{tabular}{lll}
		\toprule
		Characteristic & Group 1 & Group 2 \\
		\midrule
		Example        & 0       & 0       \\
		\bottomrule
	\end{tabular}
\end{table}

% Example figure placeholder:
%
% \begin{figure}[ht]
% \centering
% \includegraphics[width=0.8\textwidth]{figure1.pdf}
% \caption{Figure title and complete legend.}
% \end{figure}

% ==================================================================
% SUPPLEMENTAL MATERIAL
% ==================================================================

\clearpage
\section*{Supplemental Material}

 % If a survey was used, JGME requires it as supplemental material.
 %
 % Supplemental files:
 % - Appendix 1: Survey instrument
 % - Appendix 2: Interview guide
 % - Appendix 3: Intervention materials
 % - Appendix 4: Prompt/protocol library for AI intervention
 % - Appendix 5: Analytic code or codebook
 %
 % Cite each supplemental item in the main text.

 [LIST AND DESCRIBE SUPPLEMENTAL MATERIAL.]

% ==================================================================
% FINAL PRE-SUBMISSION CHECKLIST
% ==================================================================

\clearpage
\section*{JGME Pre-Submission Checklist}

\subsection*{Category and Limits}

\begin{itemize}[label=$\square$, leftmargin=*]
	\item Correct JGME article category selected.
	\item Main-text word count is within the category limit.
	\item Abstract word count is within the category limit.
	\item Number of tables/figures/boxes is within the category limit.
	\item Any justified word-limit excess is explained in the cover letter.
\end{itemize}

\subsection*{Title Page}

\begin{itemize}[label=$\square$, leftmargin=*]
	\item Every author's degree(s) are listed.
	\item Every author's affiliation and role are listed.
	\item Corresponding author is clearly identified.
	\item Business address is included.
	\item Phone number is included.
	\item Preferred email is included.
	\item Manuscript word count is included.
	\item Abstract word count is included when applicable.
\end{itemize}

\subsection*{Manuscript Formatting}

\begin{itemize}[label=$\square$, leftmargin=*]
	\item Pages are double-spaced.
	\item Pages are numbered.
	\item Abbreviations are defined and a list is supplied when used.
	\item References are numbered in order of first citation.
	\item Reference numbering is consistent in text and graphics.
	\item Every table, figure, box, and supplement is cited in the text.
	\item Tables/figures/boxes have complete legends.
	\item Footnotes use superscripted letters.
\end{itemize}

\subsection*{Research Manuscripts}

\begin{itemize}[label=$\square$, leftmargin=*]
	\item Structured abstract uses Background, Objective, Methods, Results,
	      Conclusions.
	\item Main manuscript uses Introduction, Methods, Results, Discussion,
	      Conclusions.
	\item Study year appears in abstract Methods.
	\item Study year appears in manuscript Methods.
	\item Response rate, when applicable, includes eligible N, actual N, and
	      percentage in both abstract and manuscript Results.
	\item IRB/ethics approval or exemption appears at the end of Methods.
	\item Discussion contains an explicit limitations section/statement.
	\item Survey instrument is supplied as supplemental material when used.
\end{itemize}

\subsection*{Figures}

\begin{itemize}[label=$\square$, leftmargin=*]
	\item Figure files are supplied in original format.
	\item Figures are grayscale.
	\item Figures are at least 200 dpi.
\end{itemize}

\subsection*{AI-Specific Review}

\begin{itemize}[label=$\square$, leftmargin=*]
	\item AI tool/model and role are described when relevant.
	\item Human review/verification is described.
	\item Authors retain responsibility for final content.
	\item AI use in manuscript preparation is disclosed when applicable.
	\item PHI/privacy handling is described when relevant.
	\item Institutional approval/access restrictions are described when relevant.
	\item Bias/equity implications are considered when relevant.
	\item Effects on learner reasoning/skill development are considered.
	\item AI-detection software is not treated as definitive proof of AI use.
\end{itemize}

\subsection*{Cover Letter}

\begin{itemize}[label=$\square$, leftmargin=*]
	\item Manuscript title and article category are named.
	\item Relevance to GME/JGME readership is explained.
	\item Novel contribution is summarized.
	\item Originality/concurrent-submission status is disclosed.
	\item Authorship/accountability is confirmed.
	\item Ethics status is stated.
	\item Funding/conflicts are stated.
	\item Related abstracts/preprints/presentations are disclosed.
	\item Any word-limit exception is explained.
	\item AI Teaching Rounds relevance is mentioned only if appropriate.
\end{itemize}

\end{document}
